\chapter{Conclusiones y Plan de Trabajo}
\label{sec:conclusiones}

El anteproyecto presentado demuestra la viabilidad técnica y metodológica de construir una máquina de control numérico
computarizado de bajo costo reutilizando componentes electromecánicos de impresoras fuera de uso, al tiempo que integra
de manera armónica los contenidos teóricos y prácticos de tres asignaturas troncales del cuarto año de Ingeniería
Electrónica.

\section{Resumen de Integración Técnica}
\label{sec:resumen_integracion}

La tabla~\ref{tab.ej5:resumen_materias} sintetiza la distribución de responsabilidades técnicas e hitos de diseño por
asignatura.

\begin{table}[!ht]
  \centering
  \resizebox{\textwidth}{!}{
    \begin{tabular}{|l|p{6cm}|p{6cm}|}
      \hline
      \textbf{Asignatura} & \textbf{Módulo del Proyecto} & \textbf{Conceptos Clave Aplicados} \\ \hline
      Electrónica Aplicada II & Sensado de corriente y final de carrera electrónico & Amplificador diferencial, amplificador operacional, filtro RC, comparador con histéresis (Schmitt trigger), acondicionamiento analógico. \\ \hline
      Técnicas Digitales II & Controlador central embebido e interfaz HMI & RTOS (FreeRTOS), interpretador de código G, buses SPI e I2C, conversor ADC, interrupciones externas EXTI, timers por hardware. \\ \hline
      Máquinas e Inst. Eléctricas & Controladores modulares de motores PaP y Spindle & Motores paso a paso bipolares, control modular por canal, modulación PWM, puentes H, diodo flyback, evaluación de micropasos. \\ \hline
    \end{tabular}
  }
  \caption{Matriz de integración conceptual por asignatura académica.}
  \label{tab.ej5:resumen_materias}
\end{table}

\section{Pasos Futuros para la Etapa de Desarrollo Final}
\label{sec:pasos_futuros}

Con la aprobación del presente anteproyecto, el equipo de trabajo avanzará en las siguientes etapas del proyecto final:
\begin{enumerate}
  \item Caracterización mecánica del chasis de impresora recuperado y selección final de motores.
  \item Simulación analógica en LTSpice de la etapa de sensado con OpAmps y comparador de Electrónica Aplicada II.
  \item Diseño de los esquemáticos y trazado de placas de circuito impreso (PCB) en KiCad para los controladores modulares
    de potencia.
  \item Desarrollo del firmware embebido sobre FreeRTOS en lenguaje C/C++ verificando tiempos de ejecución y latencias.
  \item Ensamble, calibración y pruebas de mecanizado real sobre placas de circuito impreso y materiales de prueba.
\end{enumerate}
